\documentclass[12pt, letterpaper]{article}

% ─── Paquetes esenciales ───────────────────────────────────────────────────────
\usepackage[utf8]{inputenc}
\usepackage[T1]{fontenc}
\usepackage[spanish]{babel}
\usepackage{geometry}
    \geometry{margin=2cm}
\usepackage{graphicx}
\usepackage{amsmath, amssymb, amsthm}
\usepackage{physics}
\usepackage{xcolor}
\usepackage{booktabs}
\usepackage{hyperref}
    \hypersetup{
        colorlinks = true,
        linkcolor  = udea,
        urlcolor   = udea,
        citecolor  = udea
    }
\usepackage{fancyhdr}
\usepackage{titlesec}
\usepackage{enumerate}
\usepackage{lipsum}

% ─── Color institucional UdeA ─────────────────────────────────────────────────
\definecolor{udea}{RGB}{0, 110, 55}    % verde institucional UdeA

% ─── Datos del documento ──────────────────────────────────────────────────────
\newcommand{\curso}{Mecánica Celeste}
\newcommand{\tarea}{Problem Set N.º\,1}
\newcommand{\periodo}{2026-1}

\newcommand{\autores}{%
    Sol\\[4pt]
    Sara\\[4pt]
    Daniel%
}

\newcommand{\institucion}{Universidad de Antioquia}
\newcommand{\facultad}{Facultad de Ciencias Exactas y Naturales}
\newcommand{\departamento}{Instituto de Física}

% ─── Encabezado y pie de página ───────────────────────────────────────────────
\pagestyle{fancy}
\fancyhf{}
\fancyhead[L]{\small{\curso\ –\ \tarea}}
\fancyhead[R]{\small{\periodo}}
\fancyfoot[C]{\thepage}
\renewcommand{\headrulewidth}{0.4pt}

% ══════════════════════════════════════════════════════════════════════════════
\begin{document}
\thispagestyle{empty}

% ─── PORTADA ─────────────────────────────────────────────────────────────────
\begin{titlepage}
    \centering

    % Logo
    \includegraphics[width=0.28\textwidth]{logoUdea.png}\\[0.8cm]

    % Línea decorativa superior
    {\color{udea}\rule{\linewidth}{1.5pt}}\\[0.4cm]

    % Institución y dependencia
   % {\large\bfseries\color{udea} \institucion}\\[0.15cm]
    {\normalsize \facultad}\\[0.05cm]
    {\normalsize \departamento}\\[0.6cm]

    % Título principal
    {\LARGE\bfseries \tarea}\\[0.3cm]
    {\Large \curso}\\[0.2cm]
    {\normalsize Período académico: \periodo}\\[0.5cm]

    {\color{udea}\rule{\linewidth}{0.8pt}}\\[0.8cm]

    % Bloque de autores (centrado)
    {\small\bfseries Autores}\\[0.4cm]
    {\small \autores}

    \vfill

    % Pie de portada
    {\color{udea}\rule{\linewidth}{1.5pt}}\\[0.2cm]
    {\footnotesize Medellín, Colombia — \the\year}

\end{titlepage}

% ─── TABLA DE CONTENIDOS ─────────────────────────────────────────────────────
\newpage
\thispagestyle{empty}           % sin encabezado en la página del índice
{
    \hypersetup{linkcolor=black} % los enlaces del índice en negro, más limpio
    \tableofcontents
}
\newpage                        % el cuerpo empieza en página nueva

% ─── CUERPO DEL DOCUMENTO ────────────────────────────────────────────────────
\section{Problema 1}

Sea el vector de posición definido en coordenadas cartesianas como $\vec{r} = (x, y, z)$ y su
magnitud escalar definida como $r = |\vec{r}| = (x^2 + y^2 + z^2)^{1/2}$.

\begin{enumerate}[a.]
    \item Calcule las derivadas parciales $\frac{\partial r}{\partial x}$,
          $\frac{\partial r}{\partial y}$ y $\frac{\partial r}{\partial z}$.
    \item A partir de lo anterior, demuestre que el gradiente de la magnitud del vector de
          posición es $\nabla r = \frac{\vec{r}}{r}$ y verifique que la magnitud de este
          gradiente es $|\nabla r| = 1$ para todo $r \neq 0$.
\end{enumerate}

% ── Solución ──────────────────────────────────────────────────────────────────
\subsection{Solución}

\subsubsection*{Parte (a): Cálculo de las derivadas parciales}

Tenemos: $r(x,y,z) = (x^2+y^2+z^2)^{1/2}$ . Derivando respecto a $x$:

$$
\frac{\partial r}{\partial x}
    = \frac{1}{2}(x^2+y^2+z^2)^{-1/2}\cdot 2x
    = \frac{x}{(x^2+y^2+z^2)^{1/2}}
    = \frac{x}{r}
$$

Por la simetría esférica de $r$, el procedimiento es el mismo para las demás
coordenadas:

$$
\frac{\partial r}{\partial y} = \frac{y}{r}, \qquad
\frac{\partial r}{\partial z} = \frac{z}{r}
$$

\subsubsection*{Parte (b): Gradiente y verificación de magnitud}

Veamos el operador gradiente

$$
\nabla r
    = \left(\frac{\partial r}{\partial x},\,
            \frac{\partial r}{\partial y},\,
            \frac{\partial r}{\partial z}\right)
    = \left(\frac{x}{r},\frac{y}{r},\frac{z}{r}\right)
    = \frac{1}{r}(x,y,z)
    = \frac{\vec{r}}{r} \equiv \hat{r}
$$

donde $\hat{r}$ es el vector unitario radial.

Ahora verificamos la magnitud

$$
|\nabla r|
    = \sqrt{\frac{x^2}{r^2}+\frac{y^2}{r^2}+\frac{z^2}{r^2}}
    = \sqrt{\frac{x^2+y^2+z^2}{r^2}}
    = \sqrt{\frac{r^2}{r^2}}
    = 1, \qquad r \neq 0
$$

% ── Relevancia ────────────────────────────────────────────────────────────────
\subsection{Relevancia en Mecánica Celeste}

\begin{enumerate}
    \item \textbf{Fuerzas centrales.} Muchas fuerzas en Mecánica Celeste y Electrodinámica
          dependen únicamente de $r$. Este resultado permite escribir
          $\vec{F} = -\nabla V(r) = -\hat{r}\,\mathrm{d}V/\mathrm{d}r$, lo que implica que
          las fuerzas gravitacionales siempre apuntan a lo largo de la línea que une a los
          cuerpos.
          
    \item \textbf{Campos conservativos.} Es la base para derivar campos vectoriales a partir
          de potenciales escalares, lo que simplifica la resolución de problemas de dinámica
          de $N$ cuerpos.
\end{enumerate}






\end{document}